\documentclass[11pt,reqno]{amsart}

\usepackage[T1]{fontenc}
\usepackage{newtxtext,newtxmath}
\let\Bbbk\relax
\usepackage{amsmath,amssymb,amsfonts}
\usepackage{microtype}
\usepackage{xcolor}
\usepackage[numbers]{natbib}
\usepackage[colorlinks=true,linkcolor=blue,citecolor=blue!60!black,urlcolor=blue!60!black]{hyperref}

\newcommand{\sg}{\operatorname{sg}}

\begin{document}

\title{Spontaneous differentiation in collective learning}
\maketitle

\section{Literature}

Collective learning can be modeled as a population of neural networks whose parameters evolve under local learning and diffusive coupling. Arola-Fernández and Lacasa introduced a minimal dynamics of this form and derived a coarse-grained disordered Ginzburg--Landau theory in which coupling competes with heterogeneity generated by the private data available to each learner \cite{arola2024}. The model exhibits a coupling-driven collective learning phase, but the heterogeneity is fixed by the data partition.

Evolutionary dynamics provides a mechanism by which heterogeneity can change with the state of the population. Selection and mutation modify the distribution of heritable strategies according to their population-dependent fitness \cite{page2002}. We use this mechanism to promote heterogeneity from a fixed input to a dynamical degree of freedom.

We consider coupled neural learners exposed to the same data distribution. Neural parameters change rapidly by backpropagation, while a slower learning strategy changes through selection and mutation. The slow variable does not define a private dataset; it changes how examples from the common data stream contribute to the local gradient. The resulting competition is between diffusive coupling, which homogenizes neural parameters, and selection, which can maintain or amplify differences in learning strategy.

\section{Results}

\subsection{Coupled learners in a common environment}

Consider $N$ neural networks with identical architecture. Learner $i$ has parameters $\boldsymbol{\theta}_i = (\theta_i^1,\ldots,\theta_i^P)$ and prediction $f(x;\boldsymbol{\theta}_i)$. At each step, all learners receive the same minibatch $B_t = \{(x_b,y_b)\}_{b=1}^{n}$ drawn from a common distribution.

Without an additional source of heterogeneity, the microscopic dynamics follows the same structure as the coupled learning model of Ref.~\cite{arola2024},
\begin{equation}
\theta_i^\alpha(t+1)
=
\theta_i^\alpha(t)
-
\frac{\eta}{n}
\sum_{b=1}^{n}
\partial_{\theta_i^\alpha}
\ell\!\left(
f(x_b;\boldsymbol{\theta}_i),
y_b
\right)
+
\frac{\eta\sigma}{N}
\sum_j q_{ij}
\left(
\theta_j^\alpha-\theta_i^\alpha
\right),
\label{eq:baseline}
\end{equation}
where $\eta$ is the learning rate, $\sigma$ is the coupling strength, and $q_{ij}$ specifies the interaction network. The first term performs local learning and the second tends to homogenize the parameters of connected learners.

If all networks see the same data and follow the same learning rule, Eq.~\eqref{eq:baseline} contains no persistent source of functional differentiation. We introduce a slow variable that changes the local learning dynamics without changing the data available to each learner.

\subsection{A slow learning genotype}

Each learner carries a learning genotype $\mathbf g_i \in \mathbb R^r$. The neural parameters $\boldsymbol{\theta}_i$ form the fast phenotype, while the genotype changes only on the slow timescale and is not optimized by backpropagation.

Let
\begin{equation}
\mathbf h_i(x)
=
\mathbf h(x;\boldsymbol{\theta}_i)
\in
\mathbb R^r
\end{equation}
denote a representation of the input generated by the current network. The genotype assigns a score
\begin{equation}
s_i(x)
=
\frac{
\mathbf g_i^{\mathsf T}
\sg[\mathbf h_i(x)]
}{T},
\label{eq:score}
\end{equation}
where $T>0$ controls the strength of the preference and $\sg[\cdot]$ denotes a stop-gradient operation. The representation can therefore change during learning, but the allocation weights are treated as fixed when differentiating the local loss with respect to $\boldsymbol{\theta}_i$.

For a minibatch $B_t$, define the normalized sample weights
\begin{equation}
w_i(x_b\mid B_t)
=
n\,
\frac{
\exp[s_i(x_b)]
}{
\sum_{c=1}^{n}
\exp[s_i(x_c)]
}.
\label{eq:weights}
\end{equation}
The normalization
\begin{equation}
\frac{1}{n}
\sum_{b=1}^{n}
w_i(x_b\mid B_t)
=
1
\label{eq:budget}
\end{equation}
fixes the total gradient budget of every learner. All samples remain available, while the genotype redistributes learning effort within the current batch.

The neural weights then evolve according to
\begin{align}
\theta_i^\alpha(t+1)
&=
\theta_i^\alpha(t)
-
\frac{\eta}{n}
\sum_{b=1}^{n}
w_i(x_b\mid B_t)
\partial_{\theta_i^\alpha}
\ell\!\left(
f(x_b;\boldsymbol{\theta}_i),
y_b
\right)
\nonumber\\
&\quad
+
\frac{\eta\sigma}{N}
\sum_j q_{ij}
\left(
\theta_j^\alpha-\theta_i^\alpha
\right).
\label{eq:microscopic}
\end{align}
Equation~\eqref{eq:microscopic} modifies only the local gradient contribution. The genotype changes the relative weight of samples in backpropagation, while diffusive coupling continues to act directly on the neural parameters.

The uniform-learning limit is recovered for $T\to\infty$, where $s_i(x)\to0$ and $w_i(x_b\mid B_t)\to1$. Equation~\eqref{eq:microscopic} then reduces to Eq.~\eqref{eq:baseline}. Finite $T$ therefore controls the genotype-dependent redistribution of learning effort.

\subsection{Selection acts on the learning strategy}

The genotype distribution evolves on the slow time $\tau = \varepsilon t$, with $0<\varepsilon\ll1$. For a large population, let $P(\mathbf g,\tau)$ denote the density of learning genotypes. A minimal mutation--selection dynamics is
\begin{equation}
\partial_\tau P
=
\left[
F(\mathbf g;P)-\overline F
\right]
P
+
\mu
\nabla_{\mathbf g}^2 P,
\label{eq:replicator}
\end{equation}
with
\begin{equation}
\overline F
=
\int
F(\mathbf g;P)
P(\mathbf g,\tau)
\,\mathrm d\mathbf g.
\end{equation}
The first term amplifies learning strategies with above-average fitness and the second represents small mutations in genotype space. In a finite population, Eq.~\eqref{eq:replicator} can be implemented by mutation and fitness-dependent resampling on the slow timescale.

We take fitness to measure contribution to collective performance rather than individual accuracy. A concrete choice is the marginal contribution of learner $i$. Define the collective predictor
\begin{equation}
\overline f(x)
=
\frac{1}{N}
\sum_i
f(x;\boldsymbol{\theta}_i)
\end{equation}
and its validation loss $L$. Let $L_{-i}$ be the same loss when learner $i$ is removed from the ensemble. Then
\begin{equation}
F_i
=
L_{-i}-L.
\label{eq:fitness}
\end{equation}
A learner has positive fitness when its removal worsens the collective prediction. The fitness is therefore frequency dependent: strategies that provide information already supplied by many learners have a smaller marginal contribution than strategies that complement the current population.

The fast and slow dynamics are coupled through fitness and sample weighting. The neural parameters determine the collective contribution $F_i$ in Eq.~\eqref{eq:fitness}; selection changes $P(\mathbf g,\tau)$ through Eq.~\eqref{eq:replicator}; and the resulting genotype distribution changes the weights $w_i$ that enter subsequent updates of Eq.~\eqref{eq:microscopic}. The parameter $\varepsilon$ controls the separation between these two timescales.

\subsection{Coarse-grained description}

The microscopic dynamics motivates an effective description in which the fixed disorder of the original collective learning problem is replaced by an adaptive field. Let $m_{i\mu}$ denote the amplitude learned by unit $i$ along a coarse-grained learning direction $\mu=1,\ldots,K$. Let
\begin{equation}
\mathbf a_i
=
(a_{i1},\ldots,a_{iK}),
\qquad
a_{i\mu}\geq0,
\qquad
\sum_{\mu=1}^{K}a_{i\mu}=1,
\label{eq:simplex}
\end{equation}
be the coarse-grained allocation generated by the genotype.

A minimal effective ansatz is
\begin{align}
\dot m_{i\mu}
&=
a_{i\mu}
R_\mu(m_{i\mu})
-
\gamma m_{i\mu}
+
\frac{\sigma}{N}
\sum_j q_{ij}
\left(
m_{j\mu}-m_{i\mu}
\right),
\label{eq:effective-m}
\\
\dot a_{i\mu}
&=
\varepsilon
a_{i\mu}
\left(
F_{i\mu}
-
\sum_\nu
a_{i\nu}F_{i\nu}
\right)
+
\varepsilon\mu
\left(
\frac{1}{K}-a_{i\mu}
\right).
\label{eq:effective-a}
\end{align}
Here $R_\mu$ denotes the effective local learning dynamics along direction $\mu$. Equations~\eqref{eq:effective-m} and \eqref{eq:effective-a} are an effective ansatz rather than a derivation from Eq.~\eqref{eq:microscopic}. They isolate the structural change introduced by the genotype: the learning rate along each coarse-grained direction becomes adaptive and evolves on the slow timescale.

The permutation-symmetric state has $a_{i\mu}=1/K$ for all $i$ and $\mu$. Diffusive coupling in Eq.~\eqref{eq:effective-m} suppresses differences between learners, while frequency-dependent fitness in Eq.~\eqref{eq:effective-a} can amplify differences in learning allocation. Its stability therefore determines whether specialization can emerge from an initially symmetric population.

Two order parameters separate collective learning from specialization. Define
\begin{equation}
M^2
=
\sum_{\mu=1}^{K}
\left(
\frac{1}{N}
\sum_i m_{i\mu}
\right)^2
\label{eq:collective-order}
\end{equation}
and
\begin{equation}
S
=
\frac{1}{N}
\sum_i
\left\|
\mathbf m_i-\overline{\mathbf m}
\right\|^2,
\qquad
\overline{\mathbf m}
=
\frac{1}{N}
\sum_i\mathbf m_i.
\label{eq:specialization-order}
\end{equation}
The quantity $M$ measures collective learning, while $S$ measures differentiation between learners. A third quantity,
\begin{equation}
G
=
\frac{1}{N}
\sum_i
\left\|
\mathbf a_i-\overline{\mathbf a}
\right\|^2,
\label{eq:genotype-order}
\end{equation}
measures diversity of the slow learning strategies.

The dynamics is controlled by the coupling strength $\sigma$, the timescale ratio $\varepsilon$, mutation $\mu$, and the preference scale $T$. The first analytical problem is to determine when the symmetric state becomes unstable to perturbations in $\mathbf a_i$, and whether that instability produces a nonzero specialization order parameter $S$.

\bibliographystyle{unsrtnat}
\bibliography{refs}

\end{document}
